\section{Alur materi dan peta konsep}

Bab-bab materi inti tersusun secara berjenjang sesuai pemetaan Sub-CPMK (Gambar~\ref{fig:peta-konsep}). Alur ini membantu menempatkan setiap topik dalam konteks keseluruhan kurikulum OBE \cite{openlogic,levin_discrete}.

\begin{figure}[H]
\centering
\begin{tikzpicture}[
  node distance=0.6cm and 1cm,
  box/.style={rectangle, draw=black, rounded corners, fill=blue!10, text width=2.6cm, minimum height=0.85cm, align=center, font=\small},
  root/.style={rectangle, draw=black, rounded corners, fill=green!20, text width=3.4cm, minimum height=0.9cm, align=center, font=\small\bfseries},
  arrow/.style={->, thick, >=stealth}
]
  \node[root] (root) {Logika Matematika\\(alur kurikulum)};
  \node[box, below left=1.2cm and 0.4cm of root] (b3) {Tahap I\\Bab 3--5:\\Proposisi, tabel, inferensi};
  \node[box, below=1.2cm of root] (b4) {Tahap II\\Bab 6--8:\\Predikat dan pemodelan};
  \node[box, below right=1.2cm and 0.4cm of root] (b5) {Tahap III\\Bab 9--11:\\Boolean, K-Map, gerbang};
  \node[box, below=1.8cm of b4] (b6) {Tahap IV\\Bab 12--14:\\Bukti, induksi, validasi};
  \draw[arrow] (root) -- (b3);
  \draw[arrow] (root) -- (b4);
  \draw[arrow] (root) -- (b5);
  \draw[arrow] (b3) -- (b4);
  \draw[arrow] (b4) -- (b5);
  \draw[arrow] (b5) -- (b6);
  \draw[arrow] (b4) -- (b6);
\end{tikzpicture}
\caption{Skema tahapan materi inti dan nomor bab}
\label{fig:peta-konsep}
\end{figure}

\subsection{Tahap I (Bab 3--5): Proposisi dan inferensi}
Bab 3 dan 4 membahas proposisi, operator logika, serta tabel kebenaran dan klasifikasi formula. Bab 5 memperdalam aturan inferensi dan penarikan kesimpulan pada logika proposisional.

\subsection{Tahap II (Bab 6--8): Predikat, kuantifier, dan pemodelan}
Bab 6--8 mengembangkan logika predikat, kuantifier, serta contoh pemodelan masalah komputasi. Konsep ini menjadi landasan untuk menerjemahkan persyaratan informal ke bahasa formal.

\subsection{Tahap III (Bab 9--11): Aljabar Boolean dan rangkaian}
Bab 9--11 membahas hukum aljabar Boolean, penyederhanaan (termasuk K-Map), dan representasi ke gerbang logika sebagai jembatan ke pemahaman struktur rangkaian digital.

\subsection{Tahap IV (Bab 12--14): Pembuktian dan validasi}
Bab 12--14 menempatkan teknik pembuktian, induksi matematika, serta argumen formal untuk validasi algoritma dan verifikasi program sesuai Sub-CPMK 4.x, termasuk konsep prekondisi, postkondisi, dan loop invariant pada bagian validasi.
